% Transcription — fonds Alexandre Grothendieck, shelfmark 141, batch 1 (pages 1-15).
% Established from the facsimile archives/batches/141/batch-01.pdf (a mirror of
% https://grothendieck.umontpellier.fr/141.pdf), per the skill
% transcribe-grothendieck. The folder is fifteen pages and this is its only
% batch.
%
% Twelve of the fifteen leaves are transcribed. Pages 2, 4 and 6 are typed
% administrative notices (USTL thesis-defence announcements of November 1980)
% carrying nothing in his hand; they take no \page marker. No page carries a
% title of his; the section headings are the edition's and say so.
%
% Four runs, on four kinds of paper.
% Pages 1, 3, 5 (blue ink, large hand): Kummer theory. V a complete
% noetherian discrete valuation ring, k residue field, K fraction field;
% with char k = 0 and k algebraically closed the extensions of K are cyclic,
% hence Kummerian, K*/K*^n = Z/nZ via the valuation, Gal(Kbar/K) = T(k) =
% lim mu_n(k) = Zhat; in residue characteristic p the prime-to-p quotient
% of pi_1(K), the exact sequence 1 -> lim mu_n(Kbar) -> pi_1(K) -> pi_1(k)
% -> 1 with pi_1(k) = pi_1(V); pi_1(S, sbar) -> lim mu_n(S) for a scheme S;
% then the Kummer sequence 0 -> H^0(S, mu_n) -> ... -> H^1(S, G_m), the
% exact sequence 0 -> G_m(S)_n -> Rev(S, mu_n) -> _nPic(S) -> 0 and a |->
% Spec O_S[T]/(T^n - a).
% Pages 7-8 (small card, black ink, a blue correcting pen): Thm 1, the
% classifying space of a discrete group pi — hot(X, B) -> H^1(X, pi) and its
% bijectivity for all X, equivalent to E contractible (the card breaks off
% at « 2) E est contractile »); Thm 2, existence of a universal pi-bundle,
% and for a pi-space X the Borel construction E_pi x_pi X with its two
% projections and the homotopy exact sequence ending pi -> pi_0(X).
% Pages 9-10 (blue felt pen, sketch pages): the world of cartes — the
% cartographic groups Gamma_{p,q}, Gamma_{infty,3} and its oriented subgroup
% SL(2,Z)/(+-1) with quotient S_3, the categories C_{p,q} and the functors
% C_{p,q} -> C_{p',q'} = C_{p,q}/R for p'|p, q'|q; lattices omega_1, omega_2
% with Im(omega_2/omega_1) > 0; « carte trouée » (X ⊃ K) whose Déc(X, K) has
% components homeomorphic to standard punctured discs D*_n, 1 <= n <= infty,
% built from the regular polygon Pi_n; and a page of drawings.
% Pages 11-15 (thin grey paper, black ink; his own numbers ① and ② on 11
% and 12, none after): pi-sets and bitorsors — index formula d'-1 = n(d-1)
% for free subgroups, the triangle group pi^+ with rho_0, rho_1, rho_infty;
% functors Ens(pi) -> C commuting with inductive limits as pi-objects of C,
% F_X(E) = X x_pi E; Hom_!(Ens(pi), Ens(pi')) = (pi', pi)-sets; Proposition
% (F left exact iff X is a right pi-torsor), Corollaire (F an equivalence
% iff X is a (pi', pi)-bitorsor), Équiv(Ens(pi), Ens(pi')) = Bit(pi', pi);
% involutive bitorsors, Bit.invol(pi) = Invol(Ens(pi)) = Ext(Z/2, pi); the
% groups pi_I = (Z/2)^{*I} for n-gonal cartes; a page on the fundamental
% groupoid of the plane minus points (lambda_i, epsilon) and a product law
% (gamma', l')(gamma, l) = (gamma' gamma, l' o gamma'(l)) for pi_1(X, Gamma,
% a); composition of bitorsors F_Y o F_X = F_{Y ^_{pi'} X}, pi_0(Bit(pi',
% pi)) = Isomext(pi', pi), and an involution F^2 = id with F mu = mu F.
%
% Reading hazards fixed once here. His underlining is set in bold. « Ens »
% is set \mathrm{Ens}, « Bit » \mathrm{Bit}; the left-and-right subscripts of
% a biset are set {}_{\pi'}X_{\pi}. The twisted group he writes with a left
% superscript X over pi is set {}^{X}\pi. « hot » (homotopy classes) is kept.
% His \varinjlim is a « lim » with a right arrow under it, \varprojlim with a
% left arrow; both are set as such. Drawings are described in \note, never
% redrawn; the positions of the scattered formulas of pages 9, 10, 13 and 14
% are not reproduced, only their content, in reading order top-left to
% bottom-right.
%
% Readings flagged and not settled: « a.v.d. » on page 1 (reads « W.d. »);
% « donc » on page 1; the condition on n on page 3 (« n premier à s »); the
% superscript target of the second vertical arrow on page 5 and the « _nPic »
% on its third line; « bijectif » in blue over another word on page 7; the
% word after « opère » and the target « X/pi » on page 8; « Gamma_1 » in the
% vertical sequence and « 2^3.3^4 » on page 9; « variété », « pol. conv. »,
% « côtés » and « F_n » on page 10; « avec » on page 11; the right vertical
% arrow on page 12; « fini », « pondérées » and « L_I/- » on page 13; the
% endpoints of the three paths on page 14; « passage au quotient » and « à
% Isomext près » on page 15.
%
% Pass: Opus 5 (claude-opus-5), 2026-09-13 — first pass, unchecked against the pages by a human.
\documentclass[11pt,a4paper]{article}
% Le préambule commun à toutes les transcriptions du fonds Grothendieck.
%
% Il tient en une idée : **l'incertitude se compose, elle ne se résout pas.**
% Une transcription qui lisse un mot douteux a détruit la seule chose pour
% laquelle on la faisait — un lecteur doit pouvoir savoir, à chaque endroit,
% ce qui était sur le feuillet et ce qui a été ajouté. Les six macros qui
% suivent sont donc l'appareil critique, et non de la décoration.
%
% Les mêmes commandes sont reconnues par `scripts/render.mjs`, qui produit la
% vue de lecture du site. Une macro ajoutée ici doit l'être aussi là-bas,
% faute de quoi le rendu échouera bruyamment — ce qui est voulu : un rendu
% silencieusement faux est pire qu'un rendu absent.

\ProvidesPackage{grothendieck}[2026/08/08 Transcriptions du fonds Grothendieck]

\usepackage{amsmath,amssymb,amsthm}
\usepackage{mathrsfs}
\usepackage{stmaryrd}
% Les diagrammes commutatifs sont partout dans ce fonds : c'est la langue
% même de la géométrie algébrique de Grothendieck, et une transcription qui
% les rendrait en prose ne transcrirait rien.
\usepackage{tikz-cd}
% Français par défaut — c'est la langue des feuillets — mais l'anglais est
% chargé aussi : la traduction et le résumé sont des documents anglais, et la
% typographie française (espaces avant les deux-points) y serait une faute.
% Ces documents basculent par \selectlanguage{english} en tête de corps.
\usepackage[english,french]{babel}
\usepackage{fontspec}
\usepackage{geometry}
\geometry{a4paper,margin=25mm,marginparwidth=28mm}
\usepackage{marginnote}
\usepackage{xcolor}
% ulem, not soul. `soul`'s \st and \ul are built on a character-by-character
% scan that XeLaTeX's font machinery breaks: the document fails with an
% undefined control sequence rather than degrading. ulem does the same two jobs
% and is Unicode-engine safe. `normalem` stops it hijacking \emph, which the
% transcriptions use for Grothendieck's own emphasis.
\usepackage[normalem]{ulem}
\usepackage{hyperref}
\hypersetup{colorlinks=true,linkcolor=black,urlcolor=black,pdfborder={0 0 0}}

% --- Métadonnées du lot -----------------------------------------------------
% Reprises telles quelles par le script de rendu, qui les affiche en tête de
% la vue de lecture. Elles fixent ce que le fichier prétend couvrir : sans
% elles, une transcription flotte sans point d'attache dans un fonds de seize
% mille feuillets.

\newcommand{\folder}[1]{\gdef\grothFolder{#1}}
\newcommand{\batch}[1]{\gdef\grothBatch{#1}}
\newcommand{\leaves}[2]{\gdef\grothFirst{#1}\gdef\grothLast{#2}}
\newcommand{\foldertitle}[1]{\gdef\grothFoldertitle{#1}}
\gdef\grothFolder{?}\gdef\grothBatch{?}\gdef\grothFirst{?}\gdef\grothLast{?}\gdef\grothFoldertitle{}

% --- Le résumé --------------------------------------------------------------
%
% Il ouvre la lecture modernisée, sous le titre « Résumé », et il est composé
% en petit. Ce n'est pas un ornement : c'est le seul passage du document qui
% ne soit pas des mathématiques, et un lecteur doit pouvoir le reconnaître
% avant de l'avoir lu — puis le sauter s'il connaît déjà le sujet, ou s'y
% arrêter s'il ne le connaît pas. Le filet qui le suit marque la reprise du
% corps.

\newenvironment{resume}
  {\small\setlength{\parskip}{0.45em}}
  {\par\medskip\hrule\medskip}

% --- Mots-clés ---------------------------------------------------------------
%
% En anglais, en clôture du résumé de la lecture modernisée : le vocabulaire
% moderne sous lequel chercher ce que les feuillets construisent. Le manifeste
% du site (`npm run manifest`) les extrait pour étiqueter la cote entière —
% cette ligne est la source unique des tags, il n'y a pas de fichier à côté.
\newcommand{\keywords}[1]{\par\smallskip\noindent
  {\footnotesize\textsc{Keywords} --- \textit{#1}}\par}

% --- L'appareil critique ----------------------------------------------------

% Le numéro de feuillet, en marge. C'est l'ancre qui permet de régler un
% désaccord sur une formule en regardant *un* feuillet plutôt qu'une chemise
% de six cents. Le site s'en sert aussi pour tourner les pages du fac-similé
% quand on fait défiler la transcription.
\newcommand{\leaf}[1]{%
  \marginnote{\footnotesize\color{gray!70}\textsf{#1}}%
  \label{leaf:#1}\ignorespaces}

% Illisible : on ne devine pas. Un mot inventé qui se lit comme les autres est
% le pire résultat possible.
\newcommand{\ill}{\textcolor{red!60!black}{[\,\ldots\,]}}

% Lecture douteuse : le mot est donné, et signalé comme incertain.
\newcommand{\uncertain}[1]{\uline{#1}}

% Ajout de l'éditeur — une lettre manquante, un mot sous-entendu.
\newcommand{\add}[1]{\textcolor{blue!50!black}{[#1]}}

% Biffé par Grothendieck lui-même. Ce qu'il a rayé fait partie du document :
% une rature dit souvent où la pensée a changé de direction.
\newcommand{\struck}[1]{\sout{#1}}

% Note du transcripteur, sur l'état matériel du feuillet.
\newcommand{\note}[1]{\par\smallskip{\small\color{gray!60!black}\textit{[#1]}}\par\smallskip}

% Note marginale de Grothendieck, distincte du corps du texte.
\newcommand{\marginal}[1]{\par\smallskip\noindent\hspace{1em}%
  {\small\color{gray!60!black}#1}\par\smallskip}

% --- Titre ------------------------------------------------------------------

\AtBeginDocument{%
  \begin{center}
    {\large\bfseries Fonds Alexandre Grothendieck --- cote n\textsuperscript{o}~\grothFolder}\\[2pt]
    {\itshape\grothFoldertitle}\\[4pt]
    {\small feuillets \grothFirst--\grothLast\ (lot \grothBatch)}\\[2pt]
    {\footnotesize\color{gray!60!black}Transcription établie d'après le fac-similé numérisé
     par l'Université de Montpellier.\\
     Les lectures incertaines sont soulignées, les illisibles notées [\,\ldots\,],
     les ajouts entre crochets.}
  \end{center}
  \vspace{1em}}

\endinput


\folder{141}
\batch{1}
\pages{1}{15}
\dating{[à partir de 1980]}
\watermark{Édition de démonstration}
\foldertitle{Théorie de Kumer ? + note Alexandre (bi tenseur…) : notes manuscrites (s.d.)}

\begin{document}

\section{Théorie de Kummer sur un anneau de valuation discrète}

\note{ce titre est de l'édition : aucune page du dossier ne porte de titre de
sa main. Les pages 1, 3 et 5 sont écrites au stylo bleu, en grands caractères,
sur des feuillets dont les versos sont des pages 2, 4 et 6 sans rapport avec
les mathématiques}

\page{1}

$V$ : \textbf{\uncertain{a.v.d.} noethérien} ; complet

$\mathrm{Spec}(V)$ \qquad $k$ corps résiduel \qquad $K$ corps des fractions

\struck{Théor} Théorie galoisienne $(K)$

a) $\mathrm{car}(k) = 0$ \qquad \textbf{$k$ alg. clos}

\note{dans $\mathrm{car}(k)$, la lettre est surchargée : un $K$ écrit sur un
$k$, ou l'inverse ; la suite (« $k$ alg. clos ») veut le corps résiduel}

\uncertain{les} extens\add{ions} \uncertain{de} $K$ sont abél\add{iennes}
cycliques \uncertain{donc} kummeriennes.

\[
K^{*}/K^{*n} \simeq \mathbb{Z}/n\mathbb{Z} \qquad \text{(par la valuation)}
\]

Soit $\overline{K}$ clôture a\add{lgébrique} de $K$.

\[
\Gamma(\overline{K}, K) \overset{\text{déf}}{=} T_{\infty}(k)
= \varprojlim_{n} \underbrace{\mu_n(k)}_{\simeq\, \mathbb{Z}/n\mathbb{Z}}
\qquad (\simeq \widehat{\mathbb{Z}})
\]

\note{la page porte « $\simeq \widehat{2}$ » : le $\mathbb{Z}$ chapeau est
tracé comme un 2}

si $k$ car \uncertain{$\neq 0$}, $V$ pas néces\add{sairement} complet,

\[
\pi_1(K)^{\text{premier à } p} \longrightarrow
\varprojlim_{n \text{ premier à } p} \mu_n(k)
\]

\note{« premier à $p$ » est écrit sous $\pi_1(K)$ et sous la limite. Le signe
entre « car » et « 0 » est un rond surchargé ; « $\neq$ » est ce que demande la
présence de $p$}

\page{3}

$\bar{s} \in S$, \quad $\bar{s}$ \textbf{pt géométrique}
\[
\pi_1(S, \bar{s}) \longrightarrow \varprojlim_{n \text{ \uncertain{premier à} } s} \mu_n(S)
\]

\note{la condition sous la limite se lit « $n$ premier à $s$ » ; ce que
l'argument demande est vraisemblablement que $n$ soit inversible sur $S$, ce
que la page ne dit pas}

\begin{tikzcd}
1 \arrow[r] & \displaystyle\varprojlim_{n \text{ premier à } p} \mu_n(\overline{K}) \arrow[r] & \pi_1(K) \arrow[r] & \pi_1(k) \arrow[r] \arrow[d, no head, "=" description] & 1 \\
 & & & \pi_1(V) &
\end{tikzcd}

\note{l'égalité verticale est tracée comme deux traits parallèles sous
$\pi_1(k)$}

\page{5}

\[
0 \to H^0(S, \mu_n) \to H^0(S, \mathbb{G}_m) \to H^1(S, \mu_n) \to H^1(S, \mathbb{G}_m)
\]

\note{la ligne omet le second $H^0(S, \mathbb{G}_m)$ et la puissance $n$-ième
qui y mène ; la ligne suivante les rétablit}

\begin{tikzcd}
\mu_n(S) \arrow[r] & \mathbb{G}_m(S) \arrow[r] & \mathbb{G}_m(S) \arrow[r] & \mathrm{Rev}(S, \mu_n) \arrow[d] \\
 & & & \mathrm{Pic}(S) \arrow[d, "\otimes n"] \\
 & & & \mathrm{Pic}(S)
\end{tikzcd}

\note{l'étiquette « $\otimes n$ » est écrite au bout de la seconde flèche
verticale, contre le second $\mathrm{Pic}$, dont elle cache la dernière lettre ;
lue comme la puissance tensorielle $n$-ième}

\[
0 \to \mathbb{G}_m(S)_n \longrightarrow \mathrm{Rev}(S, \mu_n) \to
\uncertain{{}_{n}\mathrm{Pic}}(S) \to 0
\]

\note{« $\mathbb{G}_m(S)_n$ » est écrit tel quel, indice $n$ en bas ; le
dernier terme est surchargé}

\[
a \longmapsto \mathrm{Spec}\bigl(\mathcal{O}_S[T]/(T^n - a)\bigr),
\qquad A[T]/(T^n - a)
\]

\section{Espace classifiant d'un groupe discret}

\note{ce titre est de l'édition. Les pages 7 et 8 sont deux fiches de petit
format, à l'encre noire, avec quelques corrections au stylo bleu}

\page{7}

\textbf{Thm 1} Soit $\pi$ groupe discret, $E$ revêtement $\pi$-principal d'un
espace $B$. Considérons pour \uncertain{tout} espace $X$ l'application
\[
\mathrm{hot}(X, B) \overset{h_X}{\longrightarrow} H^1(X, \pi)
\qquad \Bigl(\overset{\text{déf}}{=} \text{classes d'isom. de $\pi$-rev. principaux sur $X$}\Bigr)
\]
fonctorielle en $X$.

\note{« tout » est ajouté au stylo bleu au-dessus de la ligne ; au même stylo,
$\mathrm{hot}(X, B) \to H^1(X,\pi)$ est souligné d'un trait ondulé, et un signe
en forme de S est porté au-dessous}

Conditions équivalentes

1) $h$ est un isomorphisme de foncteurs, i.e. \uncertain{bijectif} pour tout $X$

\note{« bi » est écrit au stylo bleu sur un autre début de mot, vraisemblablement
« in\add{jectif} », qui est biffé d'un trait}

(On dit alors que $(E, B, \pi)$ est un $\pi$-rev. universel, $B$ s'appelle
« l'espace classifiant » pour $\pi$).

2) $E$ est contractile

\note{la fiche s'arrête sur cette ligne, coupée au bord inférieur}

\page{8}

\textbf{Thm 2} Pour tout $\pi$ discret, il existe un $\pi$-fibré universel,
i.e. un espace classifiant (défini à isom. unique près dans la cat.
homotopique…)

Soit $X$ un espace sur lequel $\pi$ opère \ill{}, et prenons
\[
E_{\pi} \times_{\pi} X = (E_{\pi} \times X)/\pi
\]

\note{au-dessus du signe $=$, un mot biffé, vraisemblablement « déf » ; à la
fin de la ligne précédente, après « opère », quelques signes («
$\supset X^{\delta}$ » ?) qui ne se lisent pas}

\begin{tikzcd}
 & (E_{\pi} \times X)/\pi \arrow[dl] \arrow[dr] & \\
B_{\pi} = E_{\pi}/\pi & & X/\pi
\end{tikzcd}

\note{le but de droite est surchargé : un $E_{\pi}$ biffé après le $X$, puis
$\pi$ au-dessus ; lu $X/\pi$ sans certitude}

\[
\underset{\substack{\| \\ 0}}{\pi_2(B_{\pi})} \longrightarrow \pi_1(X) \to
\pi_1(X, \pi) \longrightarrow \pi \longrightarrow
\underset{\substack{\| \\ 1}}{\pi_0(X)}
\]

\section{Cartes, groupes cartographiques}

\note{ce titre est de l'édition. Les pages 9 et 10 sont des pages d'esquisses
au feutre bleu : formules dispersées, flèches et dessins. Le contenu est
transcrit dans l'ordre de lecture, sans en reproduire la disposition}

\page{9}

\[
\tilde{x} \in (\widetilde{X} \supset \widetilde{K} \supset \widetilde{S}) \simeq C_{p,q}
\qquad \mathbf{\Gamma_{p,q}} \supset \pi
\qquad \mathrm{Norm}(\pi)/\pi
\]

$x \in$ \textbf{$X \supset K \supset S$} \qquad \textbf{$p, q$}

\[
\Gamma_{\infty,\infty} \supset \pi,
\qquad \Gamma_{\infty,\infty} \to \Gamma_{\infty,3},
\qquad \Gamma_{\infty,\infty} \to \Gamma_{p,q} \supset \pi
\qquad\qquad \Gamma_{2,2}
\]

\note{les deux flèches issues de $\Gamma_{\infty,\infty}$ partent en éventail,
l'une vers $\Gamma_{\infty,3}$, l'autre vers $\Gamma_{p,q}$ ; un long trait en
équerre sépare cette partie de la page du reste}

$C_{\infty,3}$, avec une flèche descendant vers une sphère marquée de trois
points sur l'équateur.

\note{la phrase est de l'édition : c'est un dessin, que la flèche relie à
$C_{\infty,3}$}

Suite verticale :
\[
1 \to \uncertain{\Gamma_1} \to \Gamma^{0}_{\infty,3} = SL(2, \mathbb{Z})/(\pm 1)
\to \mathfrak{S}_3 \to 1
\]

\note{la suite est écrite de haut en bas, flèches descendantes ; le premier
terme ressemble à $\Gamma_1$, et ce que l'exactitude demande là est le
sous-groupe noyau de $\Gamma^0_{\infty,3} \to \mathfrak{S}_3$}

\[
V \ni e, \qquad \textbf{$\omega_1, \omega_2$}, \qquad I(\omega_2/\omega_1) > 0
\qquad\qquad \omega_1, \omega_2, \quad \omega_2/\omega_1 = \tau
\]

\note{à côté, une sphère couverte de hachures obliques, deux points marqués}

\[
\{\sigma_0, \sigma_1, \sigma_2 \mid \rho_s^{3} = 1\} = \Gamma_{\infty,3} \supset \pi
\ni \rho_s, \qquad \textbf{$\rho_s^{\,p}$}
\]

\note{le « $=$ » est vertical, entre la présentation et $\Gamma_{\infty,3}$ ;
la présentation n'énumère pas les autres relations}

\[
C_{\infty,3} \longrightarrow C_{p,3}
\]

Une ellipse marquée de deux croix, $0$ et \uncertain{$2^{3}\cdot 3^{4}$}, et
d'un point, avec une flèche sur son bord ; à côté $C_{1,1}$ ; au-dessous
\[
R_{(p,p',q,q')}, \qquad
\Gamma_{p,q} \longrightarrow \Gamma_{p',q'}, \qquad
C_{p,q} \longrightarrow C_{p',q'} = C_{p,q}/R, \qquad
p' \mid p, \quad q' \mid q.
\]

\note{la phrase de présentation est de l'édition. Si les croix marquent les
valeurs exceptionnelles de l'invariant modulaire, la seconde vaudrait
$1728 = 2^{6} \cdot 3^{3}$ ; la page porte ce qui se lit $2^{3} \cdot 3^{4}$,
et on ne corrige pas}

\[
C_{2,1} \qquad C_{1,2} \qquad C_{1,1}
\]

\page{10}

Carte trouée $(X \supset K)$ tels que
\[
X \text{ surface top.}, \qquad K \text{ sous-\uncertain{variété} fermée de dim } 1
\]

\note{les deux qualifications sont écrites sous $X$ et sous $K$, reliées par
des flèches}

les comp. de $\text{Déc}(X, K)$ sont homéomorphes à un $D_n^{*}$ (disque
troué standard, $1 \leq n \leq +\infty$)

$\Pi_n$ \uncertain{pol. conv.} \ill{} à $n$ \uncertain{côtés} \qquad
$\Pi_n = (S_n, A_n, \uncertain{F_n})$ \qquad $S_n \subset |\Pi_n|$

\[
\underbrace{|\Pi_n| \times [0,1]} \supset |\Pi_n| \supset S_n
\]

\note{le crochet ouvrant de $[0,1]$ est repassé et pourrait être un « $]$ »}

\[
D_n^{*} = |\Pi_n| \times \left]0,1\right] \smallsetminus S_n
\]

\note{la page est ensuite couverte de dessins : un pentagone découpé en
secteurs depuis son centre ; une ligne pointée prolongée des deux côtés par
« $\ldots$ » (le cas $n = \infty$) ; deux arcs portant des flèches transverses et
des segments en éventail ; un cercle divisé, un rayon épaissi, marqué
$\pm \frac{1}{n}$ ; un cercle aux arcs épaissis alternés, avec une flèche de
rotation ; une carte en forme de pentagone aux sommets cerclés, avec des arêtes
pendantes et un point cerclé au centre ; un astérisque}

\[
\underbrace{\sigma_0\, \sigma_1}\, \sigma_2 \qquad \mathbf{E} \qquad
\boxed{\text{Déc}(\widehat{X}, \uncertain{W})} \qquad
\widetilde{X} \smallsetminus F
\]

\section{$\pi$-ensembles et bitorseurs}

\note{ce titre est de l'édition. Les pages 11 à 15 sont sur un papier fin,
à l'encre noire ; les pages 11 et 12 portent en haut à droite un numéro cerclé
de sa main, ① et ②, les suivantes aucun}

\page{11}

\[
n \left\{ \begin{array}{ll} L & d \text{ gén.} \\ \cup & \\ L' & d' \text{ gén.} \end{array} \right.
\qquad (d' - 1) = n(d - 1), \qquad 2 = 2 \cdot 1
\]

\note{à côté, un bouquet de cercles dessiné}

\[
\begin{array}{lcl}
L \simeq \pi_1(X) & & 1 - d \\
\cup & & \\
L' \simeq \pi_1(X') & & 1 - d'
\end{array}
\]

\note{à droite, un triangle hachuré de sommets $0$, $1$, $\infty$}

\[
\pi^{+} \qquad \rho_0 \quad \rho_1 \quad \rho_\infty
\]
\[
\rho_0 \mapsto \rho_0^{-1}, \qquad \rho_\infty \mapsto \rho_\infty^{-1}, \qquad
\rho_1 \mapsto \rho_\infty \rho_0
\]

\note{les trois correspondances sont écrites verticalement sous
$\rho_0\,\rho_1\,\rho_\infty$ ; celle de $\rho_1$ est une grosse flèche
descendante}

\[
\pi \qquad \mathrm{Ens}(\pi) \xrightarrow[\text{aux } \varinjlim]{\;\text{commutant}\;} \mathcal{C}
\]

\note{$F$ est écrit au-dessus de la flèche, au-dessus de « commutant ». De
$\mathcal{C}$ part un trait vers « \uncertain{avec} $\varinjlim$ » et vers «
$(\simeq \mathrm{Ens}(\pi'))$ », et une flèche courbe
$\mathcal{C} \xrightarrow[\text{aux } \varinjlim]{G \text{ commutant}} \mathcal{C}'$}

\[
\updownarrow
\]

$X$ $\pi$-objet (à droite) de $\mathcal{C}$

\note{« à droite » est ajouté au-dessus de la ligne ; le $X$ est repassé}

\[
F \longmapsto F(\pi_s)
\]
\[
X \longmapsto F_X : E \longmapsto \underbrace{X \times_{\pi} E}_{\in\, \mathrm{Ob}\,\mathcal{C}}
\]
\[
\mathrm{Hom}_{\mathcal{C}}(X \times_{\pi} E, Y) \simeq \mathrm{Hom}_{\pi}(E, \mathrm{Hom}(X, Y)),
\qquad Y \in \mathrm{Ob}\,\mathcal{C}
\]

\note{le $Y$ est écrit sur un autre signe barré}

\[
E \simeq \coprod_{i \in I} \pi/H_i
\]
\[
X \times_{\pi} E = \coprod_i \underbrace{X \times_{\pi} (\pi/H_i)}_{X/H_i}
\]

\page{12}

\textbf{Ex.} $\mathcal{C} = \mathrm{Ens}(\pi')$

\[
\mathbf{Hom}_{!}\bigl(\mathrm{Ens}(\pi), \mathrm{Ens}(\pi')\bigr) \simeq {}_{\pi'}\mathrm{Ens}_{\pi}
\ni {}_{\pi'}X_{\pi}
\]

\note{sous le « ! », accolé : « commutant aux $\varinjlim$ »}

\[
\begin{array}{ccl}
X & & E \in \mathrm{Ens}(\pi) \\
\updownarrow & & \\
F : & & E \longmapsto X \times_{\pi} E
\end{array}
\]

\textbf{Proposition} Pour que $F$ soit exact à gauche, il f\add{aut} et
s\add{uffit} que $X$ soit un $\pi$-torseur à droite.

\note{la proposition est marquée d'un trait vertical dans la marge gauche}

$X$ torseur sous $\pi$
\[
{}^{X}\pi \overset{\text{déf}}{=} X \times_{\pi} \mathrm{int}(\pi) = \mathrm{Aut}_{\pi}(X)
\]
\[
\pi' \longrightarrow \mathrm{Aut}_{\pi}(X)
\]

\textbf{Corollaire} Pour que $F$ soit équivalence, il f\add{au}t et
s\add{uffit} que $\pi' \xrightarrow{\sim} \mathrm{Aut}_{\pi}(X)$, i.e. que $X$
soit un bitorseur sous $\pi', \pi$ ; donc

\begin{tikzcd}
\textbf{Équiv}\bigl(\mathrm{Ens}(\pi), \mathrm{Ens}(\pi')\bigr) \arrow[d, "u \,\mapsto\, u^{[-1]}"'] & \approx & \mathrm{Bit}(\pi', \pi) \arrow[d] \\
\textbf{Équiv}\bigl(\mathrm{Ens}(\pi'), \mathrm{Ens}(\pi)\bigr) \arrow[rr, "\approx"] & & \mathrm{Bit}(\pi, \pi')
\end{tikzcd}

\note{l'étiquette de la flèche verticale de droite ne se lit pas ; sur celle
de gauche, un signe avant « $u \mapsto u^{[-1]}$ » ne se lit pas non plus}

\page{13}

\struck{\ill{}} \qquad ${}_{\pi}X_{\pi}$ \qquad
$X \circ X \overset{\mu}{\xrightarrow{\sim}} 1$

\[
\lambda e = e \lambda \quad \forall \lambda \in \pi
\]

\[
\mathbf{Bit.invol}(\pi) \xrightarrow{\approx} \mathbf{Invol}\bigl(\mathrm{Ens}(\pi)\bigr)
\xrightarrow{\approx} \mathbf{Ext}(\mathbb{Z}/2, \pi)
\]

\struck{Ext.}

\[
1 \to \pi \longrightarrow \widetilde{\pi} \overset{p}{\longrightarrow} \mathbb{Z}/2 \to 1
\]

\note{le $\mathbb{Z}$ est repassé. Sous la suite : « $p^{-1}(1)$ », avec un
crochet renvoyant à « \uncertain{él. non neutre} »}

\note{deux longs traits fléchés séparent la suite de la page}

$I$ \struck{\ill{}} pol. conv. \uncertain{fini} d'ordre $n \geq 3$

cartes $n$-gonales $I$-\uncertain{pondérées}

$n \in \mathbf{D}_{I}$ \qquad
\struck{$\pi_I = L_I/(\lambda_0 \lambda_1 \cdots \lambda_{n-1} = 1)$ \ill{}}

\note{la ligne biffée est rayée en croix ; sous elle, d'autres essais biffés
où se lisent $\lambda_0$, $\lambda_{n-1}$ et un indice $I$}

\[
\pi_I = (\mathbb{Z}/2)^{*I}
\]

\note{un signe biffé précède la parenthèse ; à droite, une suite biffée
$\lambda_0 \lambda_{n-1}\,\lambda_{n-1} \cdots \lambda_1^{-1}$}

\[
1 \to \pi_I^{+} \longrightarrow \pi_I \to \mathbb{Z}/2 \to 1,
\qquad \pi_I^{+} \simeq \underbrace{L_I/\uncertain{-}}
\]

\note{sous $\pi_I$, un petit signe courbe indicé $I$ ne se lit pas. Le
sous-groupe d'indice $2$ d'un produit libre de $|I|$ copies de $\mathbb{Z}/2$ est
libre de rang $|I| - 1$, ce que « $L_I/-$ » peut vouloir dire. À droite, le
dessin d'une carte de quadrilatères aux sommets numérotés $0$, $1$, $2$, $3$,
certaines faces hachurées}

\page{14}

\note{à gauche, une sphère marquée de $i$, $-i$, $1$, d'un point à gauche
et de $0$ au centre, avec des lacets ; à droite, des lacets autour de $0$ et
de $1$ dans le plan, pointés en $i$ et $-i$}

\[
\pi \qquad \lambda_0 \; \lambda_1 \; \lambda_2 \; \lambda_3 \qquad \varepsilon
\]
\[
\struck{\lambda_0 \lambda_1 \lambda_2 \lambda_3 = 1} \qquad
\lambda_3 \lambda_2 \lambda_1 \lambda_0 = 1 \qquad \varepsilon^{2} = 1
\]
\[
\begin{cases}
\varepsilon \lambda_0 \varepsilon^{-1} = \lambda_0^{-1} \\
\varepsilon \lambda_1 \varepsilon^{-1} = \lambda_1^{-1}
\end{cases}
\qquad \varepsilon \lambda_i \varepsilon^{-1}
\]
\[
\mathrm{Hom}(i, -i) \qquad
\underset{\substack{\wr \\ \pi}}{\mathrm{Hom}(-i, i)} \qquad
\underset{\substack{\wr \\ \pi}}{\mathrm{Hom}(i, i)}
\qquad \lambda_1^{-1} \lambda_2^{-1} \lambda_1
\qquad \ell \circ \bar{\ell}^{-1}
\]

\[
i \xrightarrow{\;\varepsilon^{-1}\;} \bar{\imath} \qquad
i \xrightarrow{\;\lambda_i\;} \bar{\imath} \qquad
\bar{\imath} \xrightarrow{\;\varepsilon\;} i
\qquad\qquad \varepsilon \circ \bar{\lambda}_i \circ \varepsilon^{-1}
\]

\note{les extrémités des trois chemins sont petites et en partie surchargées ;
les barres sont lues sans certitude}

\[
\underset{\substack{\in \\ a}}{X} \qquad \Gamma
\]
\[
\bigl(\gamma \in \Gamma, \ \ell \in \mathrm{Hom}(\gamma a, \struck{\ill{}}\, b)\bigr)
\qquad
\bigl(\gamma' \in \Gamma, \ \ell' \in \mathrm{Hom}(\gamma' b, \struck{\ill{}}\, c)\bigr)
\]

\[
\gamma a \xrightarrow{\;\ell\;} b \qquad
\gamma' b \xrightarrow{\;\ell'\;} c
\]

\note{après chacune de ces deux flèches, un essai biffé ; une troisième ligne,
entièrement biffée, finit par « $\to \gamma\gamma' c$ »}

\[
a \xrightarrow{\;(\gamma\ell') \circ \ell\;} \gamma\gamma' c \qquad
\bigl(\gamma\gamma', (\gamma\ell') \circ \ell\bigr)
\]

\[
\pi_1(X, \Gamma, a) \longrightarrow \Gamma \qquad
\gamma'\gamma a \xrightarrow{\;\gamma'(\ell)\;} \gamma' b \xrightarrow{\;\ell'\;} c
\]

\[
(\gamma', \ell')(\gamma, \ell) = \bigl(\gamma'\gamma, \ \ell' \circ \gamma'(\ell)\bigr)
\]

\note{le second terme est écrit au-dessus d'un premier essai biffé. À droite,
séparé par un trait vertical, un schéma : « $\gamma a$ — $\cdot$ » au-dessus de
« $\gamma' a$ — $\cdot$ »}

\page{15}

\[
\mathrm{Ens}(\pi) \xrightarrow{\;F_X\;} \mathrm{Ens}(\pi') \xrightarrow{\;F_Y\;} \mathrm{Ens}(\pi'')
\qquad {}_{\pi'}X_{\pi} \qquad {}_{\pi''}Y_{\pi'}
\]
\[
F_Y \circ F_X = F_{Y \wedge_{\pi'} X}
\]

${}_{\pi'}X_{\pi}$ \quad bitorseur ; classifiés à isom. près ?

\[
\pi' \xrightarrow{\sim} {}^{X}\pi \overset{\text{non can.}}{\simeq} \pi
\]

\note{au-dessus du second signe, en interligne : « non can. » puis « mais
can. \uncertain{à} $\mathbf{Isomext}({}^{X}\pi, \pi)$ \uncertain{près} »}

Donc $X$ définit $\mathbf{Isomext}(\pi', \pi)$

\struck{Classes de bit$(\pi', \pi)$}

\[
\pi_0\bigl(\mathbf{Bit}(\pi', \pi)\bigr) \xrightarrow{\;\sim\;} \mathbf{Isomext}(\pi', \pi)
\]

\begin{tikzcd}
\pi_0\bigl(\mathbf{Bit.pt}(\pi', \pi)\bigr) \arrow[d] & \simeq & \mathrm{Isom}(\pi', \pi) \arrow[d, "\text{surj.}"] \\
\pi_0\bigl(\mathbf{Bit}(\pi', \pi)\bigr) & \simeq & \mathrm{Isom}(\pi', \pi)/\mathrm{Int}(\pi) = \mathrm{Isomext}(\pi', \pi)
\end{tikzcd}

\note{la flèche de droite est oblique sur la page, de $\mathrm{Isom}(\pi', \pi)$
vers le quotient ; à côté d'elle : « surj. », « passage
\uncertain{au quotient} par $\mathrm{Int}(\pi)$ »}

\[
F^{2} \overset{\mu}{\xrightarrow{\sim}} \mathrm{id} \qquad
F(F(E)) \xrightarrow{\;\mu_E\;} E
\]
\[
F(F(F(E))) \underset{\mu_{F(E)}}{\overset{F(\mu_E)}{\rightrightarrows}} F(E)
\]
\[
F\mu = \mu F : F^{3} \to F
\]

\note{les deux flèches $F(\mu_E)$ et $\mu_{F(E)}$ sont écrites l'une au-dessus
de l'autre sur un même trait fléché ; le $E$ de « $\mu_E$ » est repassé sur un
autre signe}

\end{document}
